% from ega2-1-fr.tex
\oldpage[II]{149}
le point générique de $X$, donc tous ses points sont fermés (dans $X$ et
\emph{a fortiori} dans $F$)~; en considérant le sous-préschéma fermé réduit de
$X$ ayant $F$ pour espace sous-jacent (\hyperref[I.5.2.1-fr]{I, 5.2.1}), il
résulte donc de (\hyperref[I.6.2.2-fr]{I, 6.2.2}) que $F$ est fini et discret.
L'adhérence dans $X$ d'une partie infinie de $X$ est donc nécessairement égale
à $X$.

Lorsque $X$ est une courbe algébrique quelconque, en appliquant
(\hyperref[II.7.4.3-fr]{7.4.3}) aux composantes irréductibles de $X$, on voit
que les seuls points non fermés de $X$ sont les points génériques de ces
composantes.

\begin{corollary}[7.4.4]
\label{II.7.4.4-fr}
Soient $X$ et $Y$ deux courbes algébriques irréductibles sur $k$, $f:X\to Y$
un $k$-morphisme. Pour que $f$ soit dominant, il faut et il suffit que
$f^{-1}(y)$ soit fini pour tout $y\in Y$.
\end{corollary}

En effet, si $f$ n'est pas dominant, $f(X)$ est nécessairement une partie
\emph{finie} de $Y$ en vertu de (\hyperref[II.7.4.3-fr]{7.4.3}), donc il n'est
pas possible que $f^{-1}(y)$ soit fini pour tout point de $Y$, sinon $X$ serait
fini, ce qui est absurde (\hyperref[II.7.4.1-fr]{7.4.1}). Inversement, si $f$
est dominant, pour tout $y\in Y$ distinct du point générique $\eta$ de $Y$,
$f^{-1}(y)$ est fermé dans $X$ puisque $\{y\}$ est fermé dans $Y$
(\hyperref[II.7.4.3-fr]{7.4.3})~; d'autre part, par hypothèse, $f^{-1}(y)$ ne
contient pas le point générique $\xi$ de $X$, donc est fini en vertu de
(\hyperref[II.7.4.3-fr]{7.4.3}). Enfin, pour voir que lorsque $f$ est dominant,
$f^{-1}(\eta)$ est fini, on note que la fibre $f^{-1}(\eta)$ est un schéma
irréductible de type fini sur $k(\eta)$, et de point générique $\xi$
(\hyperref[I.6.3.9-fr]{I, 6.3.9} et
\hyperref[I.6.4.11-fr]{6.4.11}). Comme $k(\xi)$ et $k(\eta)$ sont des
extensions de type fini de $k$, de même degré de transcendance 1, $k(\xi)$ est
nécessairement une extension de degré fini de $k(\eta)$, donc $\xi$ est fermé
dans $f^{-1}(\eta)$ (\hyperref[I.6.4.2-fr]{I, 6.4.2}), et $f^{-1}(\eta)$ est
par suite réduit au point $\xi$.

Nous verrons au chapitre III qu'un morphisme \emph{propre} $f:X\to Y$ de
préschémas noethériens, tel que $f^{-1}(y)$ soit fini pour tout $y\in Y$, est
nécessairement \emph{fini}~; il s'ensuivra donc de
(\hyperref[II.7.4.4-fr]{7.4.4}) qu'un morphisme dominant propre d'une courbe
algébrique irréductible dans une courbe algébrique est \emph{fini}.

\begin{corollary}[7.4.5]
\label{II.7.4.5-fr}
Soit $X$ une courbe algébrique sur $k$. Pour que $X$ soit régulière, il faut et
il suffit que $X$ soit normale, ou encore que les anneaux locaux de ses points
fermés soient des anneaux de valuation discrète.
\end{corollary}

Cela résulte aussitôt de la condition b) de
(\hyperref[II.7.4.1-fr]{7.4.1}) et de
(\hyperref[II.7.1.6-fr]{7.1.6}).

\begin{corollary}[7.4.6]
\label{II.7.4.6-fr}
Soient $X$ une courbe algébrique réduite, $\mathcal{A}$ une
$\mathcal{R}(X)$-Algèbre cohérente réduite~; alors la fermeture intégrale $X'$
de $X$ relativement à $\mathcal{A}$ (\hyperref[II.6.3.4-fr]{6.3.4}) est une
courbe algébrique normale, et le morphisme canonique $X'\to X$ est fini.
\end{corollary}

Le fait que $X'\to X$ soit fini résulte de
(\hyperref[II.6.3.10-fr]{6.3.10})~; $X'$ est donc un $k$-schéma algébrique~;
en outre, si $x_i$ $(1\leq i\leq n)$ sont les points génériques des composantes
irréductibles de $X$, $x'_j$ $(1\leq j\leq m)$ ceux des composantes
irréductibles de $X'$, chacun des $k(x'_j)$ est extension algébrique finie de
l'un des $k(x_i)$ (\hyperref[II.6.3.6-fr]{6.3.6}), donc a un degré de
transcendance 1 sur $k$. $X'$ est donc bien une courbe algébrique sur $k$, et
en outre on sait que $X'$ est somme d'un nombre fini de schémas intègres et
normaux (\hyperref[II.6.3.6-fr]{6.3.6} et
\hyperref[II.6.3.7-fr]{6.3.7}).

\begin{env}[7.4.7]
\label{II.7.4.7-fr}
On dit qu'une courbe algébrique $X$ sur $k$ est \emph{complète} si elle est
\emph{propre} sur $k$.
\end{env}

\begin{corollary}[7.4.8]
\label{II.7.4.8-fr}
Pour qu'une courbe algébrique réduite $X$ sur $k$ soit complète, il faut et il
suffit que sa normalisée $X'$ le soit.
\end{corollary}

